\anexo{Entrevista 3: Compradora de Mercado Libre}

Tercera de las tres entrevistas semiestructuradas descriptas en la metodología general del Anexo D. Nombre utilizado es un pseudónimo, en línea con el criterio de anonimización ya explicado (Sección~\ref{sec:viabilidad-legal}).

\emph{Entrevistada}: Lucía, 27 años. Compra en Mercado Libre 1--2 veces por semana.

\emph{Objetivo}: profundizar el hallazgo de la encuesta sobre señales sociales (47,5\%) y la percepción de repetición de resultados destacados (59,1\%).

\begin{enumerate}
	\item \textbf{¿Qué tan seguido comprás en Mercado Libre y qué tipo de productos?} Contó que compra bastante seguido, una o dos veces por semana, principalmente artículos de decoración, gadgets, cosas para la casa y regalos, muchas veces por impulso al ver algo en redes sociales.
	\item \textbf{Contame tu última compra de algo nuevo o novedoso: ¿cómo llegaste a ese producto?} Relató que su última compra novedosa fue un organizador de cables que vio en un video de TikTok: el video le mostró el producto, lo buscó en Mercado Libre y compró la publicación con más ventas y reseñas con fotos. Resumió que llegó al producto por el video, pero eligió al vendedor por las reseñas.
	\item \textbf{Cuando algo te lo recomiendan o lo ves en redes (TikTok/Instagram), ¿cuánto influye en que lo termines comprando?} Estimó que 8 de cada 10 compras de productos nuevos las vio antes en redes sociales, pero aclaró que el video solo le genera el deseo inicial: la confianza se la dan las reseñas y las fotos de otros compradores, y que las redes sociales sin reseñas no alcanzan para que compre.
	\item \textbf{¿Confiás más en una recomendación de conocidos, en las reseñas, o en lo que la plataforma te muestra primero?} Ordenó su confianza así: primero las recomendaciones de conocidos, después las reseñas con fotos y videos, y en último lugar lo que la plataforma muestra primero: esto último, dijo, casi no le genera confianza y a veces incluso le provoca rechazo porque lo percibe como publicidad.
	\item \textbf{Cuando buscás algo, ¿sentís que te aparecen siempre los mismos productos destacados? ¿Eso te ayuda o te molesta?} Confirmó que sí, que le aparecen siempre los mismos productos con ``Envíos Full'' y la etiqueta ``Mejor vendido'', incluso filtrando por precio o novedad. Reconoció que a veces eso le ayuda por rapidez, pero también le molesta, porque siente que no ve todas las opciones y que el resultado destacado no siempre es el mejor ni el más barato.
	\item \textbf{¿Alguna vez compraste algo ``que estaba de moda'' y te arrepentiste? ¿Qué señales no viste a tiempo?} Sí: mencionó el caso de unas luces LED de moda que se veían muy bien en el video, pero resultaron de mala calidad, se rompieron a los dos meses y el vendedor ya no existía. En retrospectiva, identificó las señales que no vio a tiempo: pocas reseñas y todas recientes (algo típico de un producto con reseñas infladas), fotos robadas de internet y un precio demasiado bajo comparado con el resto del mercado.
	\item \textbf{¿Comparás precios entre vendedores o vas al primero que aparece? ¿Cómo lo hacés?} Dijo que compara precios entre vendedores salvo que sea una urgencia: ordena por precio más envío, mira ventas y reputación, y elige entre los dos o tres primeros resultados. Para compras baratas (menos de 10 dólares) no le dedica tiempo a comparar y va directo al primero con buena reputación.
	\item \textbf{¿Qué te daría más confianza para descubrir productos nuevos (historial de precios, ver si viene creciendo, opiniones)?} Señaló que lo que más confianza le daría es el historial de precios, para saber si un descuento como ``40\% OFF'' es real o está calculado sobre un precio inflado. También valoraría ver que el producto viene creciendo en ventas y desde cuándo, junto con reseñas reales con fotos; con esa información se animaría a comprarle a vendedores que no conoce.
	\item \textbf{Si pudieras ver que un producto ``viene subiendo hace dos semanas'' o que su precio está estable, ¿cambiaría tu decisión?} Afirmó que sí cambiaría su decisión: que un producto ``viene subiendo hace dos semanas'' le indica que no es una publicación fantasma ni una estafa de un día, sino que otros ya lo compraron y lo validaron. Agregó que un precio estable la tranquiliza, porque le indica que no le están cobrando de más por la moda del momento.
	\item \textbf{¿Qué es lo que más te frustra hoy al comprar algo nuevo en la plataforma?} Mencionó la publicidad camuflada que se repite, la incertidumbre sobre si el precio que ve es justo, y el miedo a las reseñas falsas o a que el producto se rompa. Resumió que comprar algo nuevo es una lotería, porque nunca se sabe si se va a recibir lo que se vio en el video.
\end{enumerate}
